\documentclass[acmsmall,review,nonacm]{acmart}

\settopmatter{printacmref=false, printccs=false}

\begin{document}

\title[Review of Optimizing Confidential Deep Learning for Real-Time Systems]%
{Review of ``Optimizing Confidential Deep Learning for Real-Time Systems''}

\author{Ajoy Dey}
\affiliation{%
  \institution{Bangladesh University of Engineering and Technology}
  \country{Bangladesh}}

\begin{abstract}
This document reviews the paper ``Optimizing Confidential Deep Learning for
Real-Time Systems'' by summarizing its key findings, identifying useful ideas
and limitations, and proposing possible extensions.
\end{abstract}

\maketitle

\noindent\textbf{Date:} 13 June 2026

\section{Review by Myself}

\subsection{Summary}

This paper introduces new time-aware scheduling models for dynamic
earliest-deadline-first (EDF) and fixed-priority rate-monotonic (RM)
schedulers to preserve the confidentiality of DNN tasks by running them inside
trusted enclaves~\cite{babar2025optimizing}. A technique is proposed that
slices each DNN layer and runs the slices sequentially in the trusted enclave.
Deep Compression is also applied to reduce the DNN model size according to the
limitations of the enclave.

However, due to the extra context-switch overheads of individual layer slices,
the paper introduces a novel layer fusion technique. Layer fusion improves
real-time guarantees by grouping multiple layers of DNN workload from multiple
tasks, allowing them to fit and run concurrently within the enclaves while
maintaining timing constraints.

The authors consider a uniprocessor real-time system and implement the ideas
on a Raspberry Pi platform running a DNN-enabled trusted operating system
(OP-TEE with DarkNet-TZ). They evaluate the approach using three DNN
architectures: AlexNet-squeezed, Tiny Darknet, and YOLOv3-tiny. Compared to
the layer-wise partitioning approach, layer fusion can schedule up to 3x more
tasksets for EDF and 5x more tasksets for RM. Moreover, the layer fusion
technique reduces context switches by up to 11.12x for EDF and up to 11.06x
for RM. Therefore, layer fusion helps provide timing guarantees for
confidential deep inference in learning-enabled real-time systems.

\subsection{The Ideas I Like}

I like the ideas of layer slicing, layer fusion, and Deep Compression. Among
these, I like the layer fusion technique the most because it significantly
reduces the number of context switches. This reduction also contributes to
higher schedulability, which is very important for real-time systems.

\subsection{Limitation}

I identify the following limitations in the proposed approach.

\begin{itemize}
  \item \textbf{Rigid layer execution:} The current system assumes that all
  neural network layers must run inside the trusted execution environment
  using EDF or RM scheduling. It does not provide flexibility to execute only
  selected layers, such as the input and output layers, when full-pipeline
  confidentiality is unnecessary.

  \item \textbf{Inflexible enclave fitting:} The current design can select and
  fuse whole slices of a layer, but it cannot split an individual layer into
  partial slices when enclave capacity is limited.

  \item \textbf{Restricted workload assumptions:} The implementation assumes
  that only DNN inference tasks use the trusted execution environment. As a
  result, it does not currently handle mixed workloads or reclaim extra slack
  time from non-inference tasks sharing the enclave.

  \item \textbf{Vulnerability to schedule-based attacks:} The current setup
  relies on the underlying TrustZone architecture without adding scheduling
  noise. This may leave the system predictable and exposed to information
  leakage from real-time schedule-based side-channel attacks.
\end{itemize}

\subsection{As an Author How I Extend}

If I were the author of the paper, I would extend the proposed ideas in the
following directions.

\begin{itemize}
  \item \textbf{Partial confidentiality execution:} I would implement a
  flexible framework that executes only critical layers, such as input and
  output layers, inside the trusted execution environment while running
  intermediate layers outside it to improve real-time performance.

  \item \textbf{Fine-grained sub-layer partitioning:} I would develop a
  sub-layer splitting algorithm that breaks whole network slices into partial
  layers. This could help the scheduler fit DNN workloads into tightly
  constrained enclave capacities and improve schedulability.

  \item \textbf{Mixed-workload slack reclamation:} I would extend the
  scheduling model to include non-inference tasks in the trusted execution
  environment and use a dynamic slack reclamation mechanism to share enclave
  availability efficiently.

  \item \textbf{Obfuscated scheduling architecture:} I would introduce
  randomized scheduling noise into the task fusion process, together with a
  dedicated schedulability analysis, to reduce the risk of schedule-based
  side-channel attacks.

\end{itemize}

\section{Review by ChatGPT}

\subsection{Summary}

This paper presents DeepTrust, a framework for enabling confidential deep
neural network (DNN) inference in real-time cyber-physical systems using
trusted execution environments (TEEs) such as ARM TrustZone. The authors
address the challenge of executing large DNN models within memory-constrained
enclaves while still meeting strict real-time deadlines. They first introduce a
layer-wise partitioning approach that executes DNN layers sequentially inside
the enclave and develop schedulability analyses for both earliest deadline
first (EDF) and rate monotonic (RM) scheduling. However, they show that
frequent transitions between the normal and secure worlds introduce substantial
context-switch overheads, significantly degrading schedulability. To overcome
this limitation, the paper proposes a novel layer fusion technique that groups
multiple DNN layers, potentially from different tasks, into a single enclave
execution, reducing the number of secure-world transitions while maximizing
enclave utilization. The approach is evaluated on a Raspberry Pi platform
running OP-TEE with three realistic DNN models: AlexNet-squeezed, Tiny
Darknet, and YOLOv3-tiny. Experimental results demonstrate that layer fusion
can schedule up to 3x more tasksets under EDF and 5x more tasksets under RM,
while reducing context-switch overheads by up to 11x compared to layer-wise
execution. These findings show that DeepTrust effectively balances security,
memory constraints, and real-time requirements, making confidential DNN
inference practical for safety-critical embedded and autonomous
systems~\cite{babar2025optimizing}.

\subsection{Ideas, Tools, Methodologies, and Limitations}

One aspect I particularly liked in this paper is its practical integration of
security and real-time systems, two areas that are often studied separately.
The authors not only recognize the memory limitations of trusted execution
environments (TEEs) but also develop a complete schedulability analysis for
both EDF and RM scheduling, making the solution applicable to real-time
cyber-physical systems. The proposed layer fusion technique is especially
elegant because it addresses the root cause of performance degradation:
frequent enclave transitions. It intelligently groups layers from multiple
tasks to reduce context-switch overhead while maximizing enclave utilization.
The experimental evaluation on a real ARM TrustZone platform further
strengthens the work's practicality. However, the approach has several
limitations. It assumes a uniprocessor system, whereas many modern embedded AI
platforms use multicore architectures. The security model primarily protects
model confidentiality and does not analyze potential side-channel attacks on
TEEs, which are a known concern in TrustZone-based systems. The framework also
relies on offline schedulability analysis and assumes relatively predictable
workloads, which may limit its applicability to highly dynamic environments.
Additionally, model compression and layer fusion may introduce implementation
complexity and could become less effective for extremely large foundation
models or DNNs with irregular architectures. Finally, the evaluation is limited
to a small set of compressed CNN-based models and a Raspberry Pi platform,
leaving open questions about scalability to larger networks, heterogeneous
accelerators, and modern edge AI hardware.

\subsection{As an Author How I Extend}

If I were the author, I would extend this work in the following concrete
directions.

\begin{enumerate}
  \item \textbf{Multicore and heterogeneous platforms:} The current framework
  assumes a uniprocessor system. A natural extension would be to develop
  scheduling and fusion algorithms for multicore processors and heterogeneous
  edge platforms that combine CPUs, GPUs, and NPUs. This would make the
  approach more applicable to modern autonomous vehicles, drones, and robotics
  systems.

  \item \textbf{Adaptive runtime layer fusion:} The proposed fusion strategy
  is largely static and based on offline schedulability analysis. I would
  investigate a runtime-adaptive fusion mechanism that dynamically adjusts
  fusion decisions based on current workload, enclave utilization, and system
  load to improve responsiveness under varying operating conditions.

  \item \textbf{Security beyond confidentiality:} While the paper focuses on
  protecting model parameters and inference execution, future work could
  address side-channel attacks, such as cache-timing, memory-access pattern
  leakage, and power analysis, against TrustZone enclaves. Integrating
  side-channel-aware scheduling and memory management would strengthen the
  overall security guarantees.

  \item \textbf{Support for modern DNN architectures:} The evaluation focuses
  primarily on compressed CNN-based models such as AlexNet-squeezed, Tiny
  Darknet, and YOLOv3-tiny. Extending the framework to transformer-based
  architectures, vision transformers, and multimodal foundation models would
  demonstrate scalability to contemporary AI workloads.

  \item \textbf{Joint optimization of compression and scheduling:} Currently,
  model compression and scheduling are treated as separate steps. A promising
  direction would be a co-design framework that jointly determines compression
  levels, layer fusion strategies, and scheduling parameters to optimize
  security, accuracy, memory usage, and deadline satisfaction simultaneously.

  \item \textbf{Energy-aware confidential inference:} Many cyber-physical
  systems operate on limited battery power. Extending the framework with
  energy-aware schedulability analysis and fusion policies could reduce both
  context-switch overhead and energy consumption while preserving timing
  guarantees.

  \item \textbf{Distributed and edge-cloud TEEs:} Future autonomous systems may
  execute parts of DNN inference across multiple trusted devices or edge-cloud
  infrastructures. Investigating distributed enclave scheduling and secure
  partitioning of DNN workloads across multiple TEEs would significantly
  broaden the applicability of the proposed approach.

  \item \textbf{Formal verification and end-to-end case studies:} Finally, I
  would validate the framework on larger real-world applications, such as
  autonomous driving perception stacks or UAV mission systems, and formally
  verify the schedulability and security properties to provide stronger
  guarantees for safety-critical deployments.
\end{enumerate}

\section{Discussion}

\subsection{Which Generative Tool Did I Use and Why?}

I used ChatGPT as the generative AI tool for this review. I used ChatGPT
because I do not have a subscription to Claude, and I usually use ChatGPT to
summarize research papers. It helped me generate a concise summary, identify
important ideas and limitations, and suggest possible future research
directions.

\subsection{Prompts Used and Prompt Refinement}

I used the following prompts.

\begin{enumerate}
  \item Please read the research paper. Write a summary of the paper,
  including key findings in one paragraph.

  \item What are the ideas/tools/methodologies you liked, and what limitations
  do you identify for the proposed approach? Write in one paragraph.

  \item If you were the author of the paper, how would you extend the proposed
  ideas? List a couple of concrete directions in points.
\end{enumerate}

Yes, I refined my prompts to get the desired results. I used Deep Research
mode in ChatGPT so that it could think for a longer time and provide more
appropriate and precise answers. I also refined my prompts to make the
responses smaller, more focused, and more precise. For example, I specifically
asked for a one-paragraph summary and a one-paragraph discussion of liked
ideas and limitations. For future directions, I asked ChatGPT to list concrete
directions in points so the answer would be easier to organize and compare
with my own understanding.

\subsection{Divergence Between My Understanding and ChatGPT's Response}

Yes, I observed a divergence between my understanding and ChatGPT's response,
especially in the future direction part. In this paper, the authors consider a
uniprocessor real-time system running on a TEE-enabled platform and use
non-preemptive enclave execution. My proposed future directions are based on
this consideration. For example, I focused on partial confidentiality
execution, fine-grained sub-layer partitioning, mixed-workload slack
reclamation, and obfuscated scheduling for schedule-based side-channel
protection.

ChatGPT's response, however, suggested extending the work toward multicore and
heterogeneous processors. I think this is a useful long-term direction, but it
moves beyond the immediate system model of the paper. Therefore, I would first
stick to my own future directions because they are more directly connected to
the paper's current assumptions and limitations. After applying those
extensions, I would then consider a multicore real-time system running on a
TEE-enabled platform. In that later stage, I would investigate both
non-preemptive and preemptive enclave execution and study how the scheduling
analysis changes in a multicore setting.

\begin{thebibliography}{1}

\bibitem{babar2025optimizing}
Mohammad Fakhruddin Babar and Monowar Hasan.
\newblock Optimizing Confidential Deep Learning for Real-Time Systems.
\newblock \emph{ACM Transactions on Cyber-Physical Systems}, 2025.

\end{thebibliography}

\end{document}
